% Part: computability
% Chapter: tm-computations
% Section: church-turing-thesis

\documentclass[../../../include/open-logic-section]{subfiles}

\begin{document}

\olfileid{tur}{mac}{ctt}
\olsection{چرچ--ٹیورنگ دا دعویٰ}

ٹیورنگ مشیناں دا مقصد
مؤثر طریقۂ کار دے تصور
دی جگہ اک درست
ریاضیاتی تعریف دینا اے۔
ٹیورنگ دا خیال سی کہ
جیڑھا بندہ مؤثر طریقۂ کار
تے ٹیورنگ مشین، دوہاں
دے تصور سمجھدا ہووے،
اوہنوں ایہ گل فطری
لگے گی کہ مؤثر
طریقۂ کار نال کیتا
جا سکّن والا ہر
کم ٹیورنگ مشین وی
کر سکدی اے۔
ایس دعوے نوں ایہ گل
وی سہارا دیندی اے
کہ مؤثر طریقۂ کار
دی جیہڑیاں ہور
درست ریاضیاتی تعریفاں
پیش کیتیاں گئیاں،
اوہ ٹیورنگ مشین
دے تصور نال
حساب ہو سکّن والے فنکشناں دے لحاظ
نال ہم ارز نکلدیاں نیں۔ یعنی اوہ بالکل
اوہناں ہی فنکشناں
دا حساب کر سکدیاں
نیں۔ ایس دعوے
نوں \emph{چرچ--ٹیورنگ
دا دعویٰ} کہندے نیں۔

\begin{defn}[چرچ--ٹیورنگ دا دعویٰ]
\emph{چرچ--ٹیورنگ دا دعویٰ}
ایہ اے کہ مؤثر طریقۂ کار
نال جیہڑا وی حساب ہو سکدا اے،
اوہ ٹیورنگ مشین وی کر سکدی اے۔
\end{defn}

چرچ--ٹیورنگ دے دعوے
نوں دو طریقیاں نال
ورتیا جاندا اے۔
پہلا طریقہ سستی
دا بہانہ بن سکدا اے۔
فرض کرو ساڈے کول
کسے شے دا حساب کرن
لئی مؤثر طریقۂ کار
دا بیان ہووے، مثلاً
“فرضی کوڈ” وچ۔
فیر اسی چرچ--ٹیورنگ
دے دعوے دا حوالہ
دے کے ایہ کہہ
سکدے آں کہ کوئی
ٹیورنگ مشین وی
اوہی فنکشن دا حساب
کردی اے، بھاویں
اسی حقیقت وچ اوہ
مشین بنائی نہ ہووے۔

چرچ--ٹیورنگ دے دعوے
دی دوجی ورتوں فلسفیانہ
طور اُتے زیادہ دلچسپ
اے۔ ایہ ثابت کیتا
جا سکدا اے کہ
کجھ فنکشن ٹیورنگ
مشیناں نال حساب
نہیں ہو سکدے۔
ایس نتیجے تے
چرچ--ٹیورنگ دے
دعوے نوں ملا کے
ایہ نتیجہ نکلدا اے
کہ اوہناں دا حساب
کسے وی مؤثر طریقے نال نہیں
ہو سکدا۔ جے کوئی ایہو جیہا مؤثر
طریقۂ کار ہوندا،
تاں چرچ--ٹیورنگ
دے دعوے مطابق
اوہدے لئی ٹیورنگ
مشین وی ہوندی۔
ایس لئی جے اسی
ثابت کر دیئے کہ
کسے فنکشن دا حساب
کرن والی کوئی
ٹیورنگ مشین نہیں،
تاں اوہدے لئی
مؤثر طریقۂ کار
وی نہیں ہو سکدا۔
خاص طور اُتے،
چرچ--ٹیورنگ دے
دعوے نوں ایہ کہن
لئی ورتیا جاندا اے
کہ رُکّن دے مسئلے دا حل
کسے ٹیورنگ مشین کول نہیں،
تے ایہ مسئلہ کسے وی مؤثر
طریقے نال حل نہیں ہو سکدا۔

\end{document}
